\documentclass[11pt]{article}

\usepackage[T1]{fontenc}
\usepackage{lmodern}
\usepackage[margin=1in,headheight=15pt]{geometry}
\usepackage{microtype}
\usepackage{longtable}
\usepackage{booktabs}
\usepackage{array}
\usepackage{enumitem}
\usepackage[table]{xcolor}
\usepackage{listings}
\usepackage{titlesec}
\usepackage{fancyhdr}
\usepackage{caption}
\usepackage{hyperref}

\definecolor{PNBlue}{HTML}{17324D}
\definecolor{PNTeal}{HTML}{0F766E}
\definecolor{PNInk}{HTML}{1F2937}
\definecolor{PNMuted}{HTML}{58657A}
\definecolor{PNPaper}{HTML}{F5F7FA}
\definecolor{PNLine}{HTML}{D7DEE8}
\definecolor{PNCode}{HTML}{F7F9FB}

\hypersetup{
  colorlinks=true,
  linkcolor=PNBlue,
  citecolor=PNTeal,
  urlcolor=PNTeal,
  pdftitle={PointNeuron1.0 Reimplementation and Evaluation on Gold166},
  pdfauthor={PointNeuron Project}
}

\titleformat{\section}
  {\Large\bfseries\color{PNBlue}}
  {\thesection}{0.75em}{}
  [\vspace{0.25em}{\color{PNLine}\titlerule}]
\titleformat{\subsection}
  {\large\bfseries\color{PNTeal}}
  {\thesubsection}{0.75em}{}
\titlespacing*{\section}{0pt}{2.2em}{0.9em}
\titlespacing*{\subsection}{0pt}{1.4em}{0.45em}

\pagestyle{fancy}
\fancyhf{}
\fancyhead[L]{\small\color{PNMuted}PointNeuron1.0 Report}
\fancyhead[R]{\small\color{PNMuted}Gold166 Evaluation}
\fancyfoot[C]{\small\color{PNMuted}\thepage}
\renewcommand{\headrulewidth}{0.4pt}
\renewcommand{\headrule}{\hbox to\headwidth{\color{PNLine}\leaders\hrule height \headrulewidth\hfill}}

\captionsetup{font=small,labelfont={bf,color=PNBlue}}
\setlist{topsep=0.35em,itemsep=0.25em}
\setlength{\parindent}{0pt}
\setlength{\parskip}{0.45em}
\renewcommand{\arraystretch}{1.15}

\lstdefinestyle{pncode}{
  basicstyle=\ttfamily\small,
  breaklines=true,
  columns=fullflexible,
  frame=single,
  rulecolor=\color{PNLine},
  backgroundcolor=\color{PNCode},
  xleftmargin=0.5em,
  xrightmargin=0.5em,
  framesep=0.55em,
  keywordstyle=\color{PNBlue}\bfseries,
  commentstyle=\color{PNMuted},
  stringstyle=\color{PNTeal}
}
\lstset{style=pncode}

\newcommand{\reportbox}[2]{%
  \vspace{0.6em}%
  \noindent\fcolorbox{PNLine}{PNPaper}{%
    \begin{minipage}{0.94\linewidth}
      \textbf{\color{PNBlue}#1}\par\vspace{0.25em}
      #2
    \end{minipage}}%
  \vspace{0.8em}%
}

\begin{document}

\begin{titlepage}
\thispagestyle{empty}
\vspace*{0.35in}

{\Large\bfseries\color{PNTeal}Technical Report}\par
\vspace{0.35in}
{\Huge\bfseries\color{PNBlue}PointNeuron1.0\\[0.15em]
Reimplementation and Evaluation}\par
\vspace{0.2in}
{\Large\color{PNMuted}Gold166 / BigNeuron Reconstruction Pipeline}\par

\vspace{0.6in}
\noindent{\color{PNLine}\rule{\linewidth}{1pt}}
\vspace{0.25in}

\begin{center}
\begin{tabular}{>{\bfseries\color{PNBlue}}p{0.28\textwidth}p{0.58\textwidth}}
Dataset & Gold166 / BigNeuron-style 3D neuron volumes \\
Promoted baseline & Guarded30 skeleton proposal with adaptive foreground-geodesic graph reconstruction \\
Main negative result & Learned connectivity GAE underperformed the geodesic initializer on the eligible validation set \\
Architectural boundary & Early hard samples fail even with oracle-derived nodes, showing a PointNeuron1.0 graph-stage limitation \\
\end{tabular}
\end{center}

\vfill
{\large PointNeuron Project}\par
{\large June 28, 2026}\par
\end{titlepage}

\begin{abstract}
This report summarizes a reimplementation and experimental evaluation of a
PointNeuron-style neuron reconstruction pipeline on Gold166 / BigNeuron-style
data. The project began as a replication of the PointNeuron paper pipeline:
raw volume processing, point-cloud construction, skeleton proposal, graph
connectivity, and final SWC reconstruction. Through iterative experiments, the
work identified a stable PointNeuron1.0 baseline based on a guarded proposal
checkpoint and an adaptive foreground-geodesic graph initializer. It also
identified two non-promoted paths: a longer proposal checkpoint that improved
validation loss but harmed end-to-end topology, and a learned connectivity GAE
that underperformed the geodesic initializer. Finally, an oracle-node diagnosis
showed that several early hard samples expose an architectural limitation of
PointNeuron1.0: the foreground-geodesic forced-tree graph stage fails even when
given GT-derived nodes. The report concludes by documenting the promoted
baseline, the known limitations, and the recommended direction for
PointNeuron2.0.
\end{abstract}

\tableofcontents
\newpage

\section*{Executive Summary}
\addcontentsline{toc}{section}{Executive Summary}

\reportbox{Current state}{
PointNeuron1.0 has a reproducible, promoted baseline for the validated stable
domain: guarded30 proposal predictions combined with adaptive
foreground-geodesic graph reconstruction. This is the path to cite as the
working implementation.
}

\reportbox{Main lesson}{
The experiments separated training loss from reconstruction quality. A longer
proposal checkpoint reduced validation loss but damaged end-to-end topology, so
checkpoint promotion now depends on topology evaluation rather than proposal
loss alone.
}

\reportbox{Conceptual finding}{
PointNeuron1.0 is best understood as a hybrid system: learned local proposal of
skeleton nodes followed by mostly rule-driven global graph construction. The
strongest results occur when both assumptions hold. The failures occur when
local node quality is not enough to rescue global connectivity.
}

\reportbox{Known limit}{
The early hard samples are not merely undertrained proposal cases. The
oracle-node diagnosis shows that the forced-tree foreground-geodesic graph
stage can fail even when the node set is already strong, which makes this a
PointNeuron1.0 architectural boundary and a natural target for PointNeuron2.0.
}

\newpage

\section{Problem Setting and Report Scope}

Neuron reconstruction is the task of converting a 3D microscopy image of a
neuron into a compact tree representation of its centerline. The input is a
volumetric image: each voxel has an intensity value, and bright connected
structures usually correspond to neuronal signal. The desired output is not
another image. It is a graph-like skeleton whose nodes trace the neuron's
branches and whose edges describe how those branches connect.

The standard file format used here is SWC. An SWC file stores a neuronal tree
as rows of node records. Each record contains a node identifier, a node type,
3D coordinates, a radius, and the identifier of its parent node. Because each
node has one parent, a valid SWC describes a rooted tree. In this report,
\emph{ground-truth SWC} means the manually curated or dataset-provided
reference reconstruction, while \emph{predicted SWC} means the reconstruction
exported by this project.

Gold166 is the local experimental dataset used for this work. It is treated as
a BigNeuron-style collection because samples are stored as raw neuron volumes
paired with SWC reconstructions and related metadata. The dataset is not
perfectly uniform: volumes differ in size, intensity distribution, anisotropy,
and annotation quality. Much of this project was therefore not only model
training, but also making the data path reliable enough that model failures
could be separated from data-loading failures.

This report is intentionally self-contained. Code paths are included so the
reader can verify or reproduce the work, but the scientific and engineering
argument should be understandable without opening the implementation.

\section{Purpose of the Project}

The purpose of this repository is to reproduce the main stages of the
PointNeuron reconstruction pipeline and test how far a PointNeuron1.0-style
system can be pushed on Gold166 data. The target reconstruction flow is:

\begin{lstlisting}
raw 3D neuron volume
-> foreground point cloud
-> point-cloud encoder
-> skeleton proposal model
-> graph initialization / connectivity
-> SWC reconstruction
\end{lstlisting}

The original PointNeuron paper motivates a learned point-cloud approach to
neuron reconstruction. The proposal document for PointNeuron2.0 then suggests
improvements such as adaptive voxel-to-point sampling, a stronger PointNeXt-like
point-cloud encoder, DETR-style skeleton proposal, and learned GATv2
connectivity. PointNeXt refers to a stronger modern point-cloud encoder family.
DETR refers to a transformer-based detection style in which a fixed set of
queries predicts objects. GATv2 refers to a graph attention network variant for
learning edge- and neighborhood-aware node representations. This project does
not fully implement PointNeuron2.0. Instead, it establishes a reproducible
PointNeuron1.0 baseline, evaluates it, and determines which failures are
fixable within the current architecture and which require a future architecture.

Conceptually, the project asks five research questions:

\begin{enumerate}
  \item Can raw Gold166 neuron volumes and SWC labels be converted into a
  reliable PointNeuron-style training and evaluation pipeline?
  \item Can a point-cloud proposal model learn useful candidate skeleton nodes
  from sparse foreground evidence?
  \item Does better proposal training loss translate into better final neuron
  reconstruction?
  \item Is a foreground-geodesic graph initializer strong enough to connect the
  proposed nodes into a correct SWC tree?
  \item When the pipeline fails, is the failure caused by weak proposal nodes or
  by a deeper limitation in the graph-connectivity architecture?
\end{enumerate}

\section{Conceptual View of the Pipeline}

The PointNeuron1.0 pipeline can be understood as a sequence of compression and
decision stages. The raw image volume contains too much information to use
directly as a graph. The reconstruction system must progressively compress that
image into a tree while preserving the biologically important structure.

\begin{center}
\begin{tabular}{p{0.24\linewidth}p{0.33\linewidth}p{0.33\linewidth}}
\toprule
Stage & Conceptual role & Main assumption \\
\midrule
Foreground extraction & Convert the image into candidate evidence. & Bright
voxels are likely to contain neuron signal. \\
Point cloud encoding & Learn local geometry from sparse 3D points. & Branch
shape can be inferred from neighborhoods of foreground points. \\
Skeleton proposal & Predict where centerline nodes should exist. & Local
appearance is enough to propose useful skeleton candidates. \\
Graph initialization & Decide which proposed nodes connect. & True branches
should have plausible foreground-geodesic paths between nearby nodes. \\
SWC export & Serialize the graph as a rooted neuron tree. & The final graph
should be connected, tree-like, and biologically plausible. \\
\bottomrule
\end{tabular}
\end{center}

This framing matters because not every stage has the same kind of difficulty.
Proposal prediction is mostly local: it asks whether a point and its neighbors
look like they belong near a centerline. Connectivity is more global: it asks
which proposed nodes belong to the same branch, where branches split, and which
nearby-looking structures should remain separate. A reconstruction can therefore
have visually reasonable proposal nodes but still fail as a graph.

\reportbox{Central conceptual distinction}{
PointNeuron1.0 learns where skeleton nodes may be, but its strongest promoted
connectivity stage is still largely rule-driven. The experiments therefore test
not only a neural proposal model, but also the boundary between learned local
geometry and hand-designed global topology.
}

\subsection{PointNeuron1.0 Versus PointNeuron2.0}

In this report, \emph{PointNeuron1.0} refers to the implemented baseline:
foreground point-cloud construction, DGCNN/EdgeConv proposal prediction, and
foreground-geodesic graph reconstruction. It is a practical reproduction and
extension of the PointNeuron-style idea, but it does not claim to solve every
case in Gold166.

\emph{PointNeuron2.0} is the proposed next architecture. Its purpose is not
merely to tune thresholds or add more command-line options. The motivation is
conceptual: if the graph stage fails even when node locations are good, then
the next system needs a stronger way to reason about edges. That points toward
learned candidate-edge scoring, richer node and edge features, and topology
constraints applied after edge evidence has been estimated.

\subsection{Why Negative Results Matter}

Several results in this report are intentionally negative. The full60 proposal
checkpoint was not promoted. The first learned connectivity GAE was not
promoted. The early hard samples were not folded into the first full
connectivity training set. These are not abandoned experiments; they are
evidence about the architecture.

The key lesson is that neuron reconstruction should be evaluated as an
end-to-end structural problem. A model can look better locally while producing
worse topology. A graph method can produce a valid tree while inventing long
bridges. A learned edge model can achieve high recall while adding too many
false connections. The report therefore treats failure modes as measurements of
where the architecture's assumptions stop matching the data.

\section{Definitions and Metrics}

\subsection{Core Terms}

The following terms are used throughout the report.

\begin{description}[leftmargin=3.6cm]
  \item[Volume] A 3D image represented as a grid of voxels. In this project the
  volumes are Vaa3D raw files. Vaa3D is a common tool ecosystem for neuron
  visualization and reconstruction.

  \item[Voxel] One cell of the 3D image grid. A voxel has integer coordinates
  and an intensity value.

  \item[Foreground] The subset of voxels considered likely to belong to the
  neuron signal. Foreground is usually selected by intensity thresholding.

  \item[Point cloud] A sparse set of sampled foreground voxels represented as
  3D points with optional features such as intensity. PointNeuron uses point
  clouds so the model does not need to process every voxel in a dense 3D grid.

  \item[SWC] The tree-format file used to represent a neuron skeleton. Each SWC
  node has coordinates, radius, and a parent pointer.

  \item[Skeleton node] A point on the centerline of the neuron. In the ground
  truth it comes from the SWC. In prediction it comes from the proposal model
  or from an oracle diagnostic.

  \item[Proposal] A candidate skeleton node predicted from the point cloud.
  Proposals are not final reconstructions; they still need to be filtered,
  selected, connected, and exported.

  \item[Objectness] The proposal model's confidence that a sampled point is
  near a real skeleton centerline.

  \item[Center offset] A predicted displacement from a sampled foreground point
  toward the estimated skeleton center.

  \item[Radius] The estimated local radius of the neuron branch near a
  candidate skeleton node.

  \item[NMS] Non-maximum suppression. This removes redundant nearby proposals
  so that many high-confidence points on the same local structure do not all
  become graph nodes.

  \item[Checkpoint] A saved model state produced during training. A checkpoint
  can have good validation loss while still producing poor downstream
  reconstructions.

  \item[GT] Ground truth. In this report it usually means the dataset-provided
  reference SWC.

  \item[Oracle nodes] Diagnostic nodes derived from ground truth rather than
  predicted by the model. They are used to ask: if the node locations were
  already good, would the graph stage still fail?
\end{description}

\subsection{Graph Terms}

\begin{description}[leftmargin=3.6cm]
  \item[Adjacency] The edge set of a graph, often represented as a matrix where
  entry $(i,j)$ indicates whether node $i$ connects to node $j$.

  \item[Geodesic path] A shortest path measured through the foreground image
  graph rather than by straight-line Euclidean distance. In this project,
  geodesic evidence asks whether two proposed nodes can be connected by a
  plausible bright path through the original volume.

  \item[MST] Minimum spanning tree. Given candidate connection costs, an MST
  chooses a set of edges that connects all selected nodes with low total cost.
  The risk is that it always forces a tree, even when some connections are
  biologically unsupported.

  \item[Bridge edge] A fallback edge added to connect graph components when the
  constrained geodesic construction cannot connect everything normally. Bridge
  edges are useful for producing a valid single tree, but many bridges usually
  mean the reconstruction is inventing unsupported long jumps.

  \item[GAE] Graph autoencoder. Here it means a neural network that receives
  node features and an initial graph, then predicts which node pairs should be
  connected.

  \item[Top-k edge selection] Selecting the $k$ highest-scoring predicted edges,
  usually where $k$ is the number of target edges. This evaluates edge ranking
  separately from a fixed probability threshold.
\end{description}

\subsection{Success Criteria}

During the experiments it became necessary to distinguish three kinds of
success:

\begin{description}[leftmargin=3.6cm]
  \item[Proposal success] The proposal model predicts candidate nodes near the
  true skeleton.
  \item[Graph success] The selected proposal nodes are connected with edges that
  match the topology induced by the ground-truth SWC.
  \item[Reconstruction success] The graph can be converted into a valid
  single-root SWC that follows the neuron without adding unsupported bridges.
\end{description}

This distinction is important because a model can improve local proposal loss
while making the final graph worse. The main topology metrics used throughout
the project are:

\begin{description}[leftmargin=3.6cm]
  \item[Edge precision] The fraction of predicted graph edges that match the
  target graph.
  \item[Edge recall] The fraction of target graph edges recovered by the
  predicted graph.
  \item[Edge F1] The harmonic mean of edge precision and edge recall. It is
  useful because a graph with many false edges may have high recall but low
  precision, while an overly sparse graph may have high precision but low
  recall.
  \item[Bridge edges] The number of fallback edges needed to force a connected
  tree. Lower is better.
  \item[Reachable edge fraction] Fraction of final graph edges that were
  actually reachable through the foreground-geodesic image graph. Higher is
  better.
\end{description}

For evaluation, the target graph is induced from the ground-truth SWC. Predicted
or selected nodes are associated with nearby ground-truth skeleton locations,
and predicted edges are counted as correct when they agree with the local
ground-truth connectivity implied by those assignments. This makes it possible
to compare a predicted graph to the SWC tree even though the predicted nodes are
not numerically identical to the ground-truth SWC nodes.

The promotion rule used in the final stage was:

\begin{quote}
A checkpoint should not be promoted unless it improves or preserves edge F1
without creating many additional bridge edges or reducing reachable edge
fraction.
\end{quote}

\section{Repository Structure}

This section is not required to understand the experimental argument. It is a
map from the report's concepts to the repository for readers who want to inspect
the implementation or rerun the experiments.

\begin{longtable}{p{0.25\linewidth} p{0.68\linewidth}}
\toprule
Path & Purpose \\
\midrule
\texttt{configs/} & Local configuration templates, including paths such as
the Vaa3D executable. \\
\texttt{data/} & Local dataset directory. Gold166 data is expected under
\texttt{data/gold166}. \\
\texttt{docs/} & Reference PDFs, strategy notes, and this LaTeX report. \\
\texttt{scripts/} & Command-line tools for inspection, training, evaluation,
visualization, diagnostics, and reconstruction. \\
\texttt{src/pointneuron/} & Importable Python implementation of data handling,
models, and graph utilities. \\
\texttt{tests/} & Unit tests for selected model and graph behavior. \\
\texttt{tmp/} & Generated artifacts: manifests, splits, caches, checkpoints,
E2E outputs, visualizations, and reports. \\
\bottomrule
\end{longtable}

\section{Data Foundation}

The Gold166-style data is not a single clean tensor dataset. Each sample folder
can contain raw Vaa3D volumes, several SWC files, auxiliary metadata, and
occasionally invalid or misaligned labels. The first milestone was therefore a
data foundation:

\begin{enumerate}
  \item scan the dataset;
  \item select one SWC per sample;
  \item validate SWC structure;
  \item decode raw volumes;
  \item verify that SWC coordinates lie within the volume dimensions.
\end{enumerate}

SWC selection priority is:

\begin{lstlisting}
1. *swc_sorted.swc
2. stamped SWCs
3. base SWCs
\end{lstlisting}

If a preferred SWC is invalid, the scanner falls back to the next valid
candidate. This prevents downstream model errors from being caused by bad data
pairing.

Implementation pointers:

\begin{itemize}
  \item \texttt{src/pointneuron/data/gold166.py}: scans the dataset and writes
  manifests.
  \item \texttt{src/pointneuron/data/swc.py}: parses and validates SWC files.
  \item \texttt{src/pointneuron/data/vaa3d\_raw.py}: reads Vaa3D raw volumes.
  \item \texttt{src/pointneuron/data/alignment.py}: checks coordinate
  alignment.
  \item \texttt{scripts/build\_gold166\_manifest.py}: builds sample manifests.
  \item \texttt{scripts/check\_sample\_alignment.py}: verifies alignment for
  individual samples.
\end{itemize}

\section{Voxel-to-Point Transformation}

The first PointNeuron-style transformation converts the dense image volume into
a point cloud:

\begin{lstlisting}
voxels with intensity > threshold
-> (x, y, z, intensity)
\end{lstlisting}

This is intentionally close to the original PointNeuron idea: use a point-cloud
network rather than process the full voxel grid directly. The advantage is that
neuron volumes are sparse and large. The weakness is that a fixed threshold is
not equally reliable across all samples. Some volumes are noisy, saturated, or
anisotropic. These issues later become important for the early hard samples.

In practical terms, thresholding decides which voxels are allowed to become
evidence. If the threshold is too low, the point cloud includes background noise
and false structures. If it is too high, faint but real branches disappear. The
early hard samples expose this tension because their intensity distributions are
less well behaved than the stable validation range.

Implementation pointers:

\begin{itemize}
  \item \texttt{src/pointneuron/data/point\_cloud.py}: converts volumes to
  point clouds.
  \item \texttt{scripts/build\_point\_cloud.py}: writes inspectable point-cloud
  CSV files.
  \item \texttt{scripts/visualize\_sample.py}: renders foreground points and
  SWC overlays to HTML.
\end{itemize}

\section{Training Cache and Splits}

The raw volumes are too large for direct proposal training. The project builds
patch-level training records that contain foreground points and nearby SWC
skeleton nodes. These records are stored as \texttt{.npz} files and then split
deterministically into train/validation/test groups.

A patch is a smaller spatial crop sampled from a full volume. Patch training is
used because full 3D neuron volumes can contain far more foreground points than
fit comfortably into a single training example. Deterministic splits are used
so that the same records remain in the same train/validation/test groups across
runs, making checkpoint comparisons meaningful.

Each cache record contains:

\begin{lstlisting}
sampled foreground points
SWC skeleton nodes in or near the patch
edge indices and skeleton metadata
sample and patch metadata
\end{lstlisting}

Implementation pointers:

\begin{itemize}
  \item \texttt{src/pointneuron/data/training\_cache.py}: creates cached
  training records.
  \item \texttt{src/pointneuron/data/torch\_dataset.py}: loads cached records
  for PyTorch training.
  \item \texttt{src/pointneuron/data/splits.py}: builds deterministic splits.
  \item \texttt{scripts/build\_training\_cache.py}: generates training caches.
  \item \texttt{scripts/build\_split.py}: writes split files.
  \item \texttt{scripts/inspect\_dataset.py}: checks loader and tensor shapes.
\end{itemize}

\section{Encoder and Proposal Model}

The proposal stage implements a DGCNN/EdgeConv-style encoder and a proposal head
that predicts:

\begin{description}[leftmargin=3cm]
  \item[Objectness] whether a point is near the skeleton.
  \item[Center offset] how to move from the sampled point toward a skeleton
  center.
  \item[Radius] local structural radius.
\end{description}

This stage corresponds to the skeleton proposal component of PointNeuron. It
does not directly reconstruct a neuron. It produces a large set of candidate
skeleton nodes that must later be selected and connected.

DGCNN means Dynamic Graph Convolutional Neural Network. Its EdgeConv operation
builds local neighborhoods in point space and learns features from relationships
between neighboring points. This is appropriate for neuron point clouds because
branch geometry is local: nearby foreground points often carry information about
centerline direction and radius.

Implementation pointers:

\begin{itemize}
  \item \texttt{src/pointneuron/models/dgcnn.py}: DGCNN/EdgeConv encoder.
  \item \texttt{src/pointneuron/models/proposal.py}: proposal head.
  \item \texttt{src/pointneuron/models/proposal\_loss.py}: proposal losses.
  \item \texttt{scripts/train\_proposal.py}: trains proposal checkpoints.
  \item \texttt{scripts/audit\_proposal\_checkpoint.py}: audits proposal
  behavior.
  \item \texttt{scripts/aggregate\_proposals.py}: applies a checkpoint to a
  whole volume.
  \item \texttt{scripts/visualize\_proposals.py}: visualizes proposals.
\end{itemize}

\section{Graph Initialization and SWC Reconstruction}

The strongest PointNeuron1.0 graph path is the foreground-geodesic initializer.
Its basic assumption is:

\begin{quote}
Two proposal nodes should be connected if there is a plausible bright
foreground path between them through the original image volume.
\end{quote}

The graph initializer:

\begin{enumerate}
  \item filters proposal nodes by confidence and non-maximum suppression;
  \item keeps a limited graph-node budget;
  \item snaps selected nodes to foreground voxels;
  \item computes shortest paths through foreground voxels;
  \item builds an MST-like graph using geodesic distances;
  \item adds bridge edges when needed to connect components;
  \item exports edge paths to SWC.
\end{enumerate}

The phrase \emph{adaptive connected coverage NMS} refers to the promoted node
selection strategy. It begins from high-confidence proposals, suppresses nearby
duplicates, and tries to preserve spatial coverage so the selected nodes are not
all concentrated in one bright region. The word \emph{connected} reflects that
selection is judged by whether the later geodesic graph can plausibly connect
the selected nodes.

Implementation pointers:

\begin{itemize}
  \item \texttt{src/pointneuron/graph/initialization.py}: shared graph
  utilities.
  \item \texttt{scripts/initialize\_geodesic\_graph.py}: foreground-geodesic
  graph construction.
  \item \texttt{scripts/generate\_swc\_from\_graph.py}: graph-to-SWC export.
  \item \texttt{scripts/visualize\_reconstruction.py}: reconstruction HTML
  visualization.
  \item \texttt{scripts/run\_end\_to\_end.py}: full proposal-to-SWC pipeline.
  \item \texttt{scripts/evaluate\_geodesic\_baseline.py}: topology evaluation.
\end{itemize}

\section{Proposal Checkpoint Experiments}

The proposal checkpoint names encode how a model was trained. For example,
\texttt{full60} means a longer 60-epoch training run, while
\texttt{guarded30} means a 30-epoch run with more conservative losses intended
to prevent the proposal head from drifting too aggressively. These names are
project labels, not external benchmarks.

\subsection{Problem: Validation Loss Was Not Enough}

Early proposal checkpoints were evaluated with proposal-level losses and
metrics. This was insufficient. A checkpoint can reduce validation loss while
increasing center drift, increasing confidence, and harming final graph
topology.

The failed full60 checkpoint illustrates this:

\begin{lstlisting}
tmp/checkpoints/proposal_conservative_rawcenter_full60_cw3_oreg001_nw1.pt
\end{lstlisting}

It reached a better validation loss, approximately:

\begin{lstlisting}
val_loss: 5.0645
\end{lstlisting}

but harmed end-to-end topology on the hard probe:

\begin{center}
\begin{tabular}{lccc}
\toprule
Checkpoint & F1 & Bridges & Reachable \\
\midrule
old rawcenter & 0.3937 & 15.25 & 0.8878 \\
full60 & 0.4232 & 40.25 & 0.6831 \\
\bottomrule
\end{tabular}
\end{center}

Although F1 rose slightly, bridges increased dramatically and reachable fraction
collapsed. Therefore full60 was not promoted.

\reportbox{Promotion rule}{
Proposal checkpoints are not promoted by validation loss alone. A checkpoint
must preserve usable graph topology: low bridge count, high reachable fraction,
and stable reconstruction behavior after SWC export.
}

\subsection{Guarded30 Checkpoint}

The promoted proposal checkpoint is:

\begin{lstlisting}
tmp/checkpoints/proposal_conservative_guarded30_cw3_oreg005_nw3.pt
\end{lstlisting}

It was trained with conservative losses and stronger regularization:

\begin{lstlisting}
epochs: 30
batch_size: 2
lr: 0.0001
proposal_coordinate_mode: raw
loss_mode: conservative
conservative_center_weight: 3.0
offset_regularization_weight: 0.005
non_worsen_weight: 3.0
augment: True
amp: True
best_epoch: 27
best_val_loss: 5.215460332957181
\end{lstlisting}

Proposal audit:

In the audit table, \emph{proposal hit} measures how often predicted candidate
nodes land near target skeleton structure. \emph{Offset mean} is the average
size of the predicted center displacement, so a large increase can indicate
that the model is making more aggressive moves away from the sampled foreground
points. \emph{Score P90} is the 90th percentile proposal confidence; high
confidence is only useful if the resulting nodes remain topologically helpful.

\begin{center}
\begin{tabular}{lccc}
\toprule
Checkpoint & Proposal Hit & Offset Mean & Score P90 \\
\midrule
old rawcenter & 0.3897 & 0.5422 & 0.8489 \\
full60 & 0.3980 & 0.9701 & 0.9265 \\
guarded30 & 0.3910 & 0.6258 & 0.8616 \\
\bottomrule
\end{tabular}
\end{center}

Guarded30 did not dramatically improve local proposal metrics, but it preserved
controlled offsets and gave better graph behavior.

\reportbox{Promoted proposal checkpoint}{
Guarded30 is the PointNeuron1.0 proposal checkpoint because it trades a small
amount of local proposal ambition for more reliable downstream topology.
}

\section{Stable-Domain End-to-End Results}

The main stable validation range was samples 100--162. This range is called
\emph{stable} because the data path, foreground extraction, proposal selection,
and geodesic graph construction produced interpretable reconstructions without
the extreme bridge failures seen in the early hard samples. It is not claimed
to be the entire problem solved; it is the validated domain where
PointNeuron1.0 can be reported honestly.

The promoted graph configuration was:

\begin{lstlisting}
initializer: geodesic
graph selection: adaptive_connected_coverage_nms
min proposal score: 0.85
NMS mode: distance
NMS distance: 18
max graph nodes: 128
\end{lstlisting}

The label \emph{old adaptive v4} refers to the previous best adaptive geodesic
baseline before guarded30 was promoted. It is included as a control so the
reader can see whether the new proposal checkpoint improved the full
reconstruction pipeline rather than only improving an isolated training loss.

\begin{center}
\begin{tabular}{llccc}
\toprule
Range & Method & F1 & Bridges & Reachable \\
\midrule
100--109 & old adaptive v4 & 0.7822 & 1.20 & 0.9659 \\
100--109 & guarded30 & 0.7868 & 1.00 & 0.9363 \\
110--129 & old adaptive v4 & 0.7545 & 1.05 & 0.9667 \\
110--129 & guarded30 & 0.7827 & 0.65 & 0.9798 \\
130--162 & old adaptive v4 & 0.8580 & 2.1515 & 0.9308 \\
130--162 & guarded30 & 0.8708 & 2.0606 & 0.9368 \\
\bottomrule
\end{tabular}
\end{center}

Combined stable-domain comparison:

\begin{center}
\begin{tabular}{lccc}
\toprule
Method & F1 & Bridges & Reachable \\
\midrule
old adaptive v4, 100--162 & 0.8131 & 1.651 & 0.9478 \\
guarded30 adaptive v4, 100--162 & approx. 0.8295 & approx. 1.444 & approx. 0.9504 \\
\bottomrule
\end{tabular}
\end{center}

Conclusion: guarded30 is promoted for PointNeuron1.0 stable-domain proposals.

\reportbox{Stable-domain result}{
Across samples 100--162, guarded30 with adaptive geodesic reconstruction gives
the best validated PointNeuron1.0 behavior: higher F1, fewer bridges, and a
slightly higher reachable fraction than the previous adaptive baseline.
}

\section{Early Hard Samples and Architectural Diagnosis}

\subsection{Observed Failure}

The early hard samples are samples in the 0--24 range that repeatedly produced
unreliable graph topology. They are difficult because the volume geometry and
foreground signal can be unusually large, saturated, noisy, anisotropic, or
otherwise hostile to the foreground-geodesic assumption. On this range,
guarded30 improved some local behavior but graph topology became unstable:

\begin{lstlisting}
guarded30 000-024:
  samples   22
  F1        0.4709
  bridges   26.9091
  reachable 0.8049
\end{lstlisting}

The worst samples were:

\begin{lstlisting}
7, 16, 22, 23, 24
\end{lstlisting}

Score-NMS on these reused proposals made the situation worse:

\begin{lstlisting}
score_nms early worst:
  F1        0.3075
  bridges   115.4
  reachable 0.1134
\end{lstlisting}

This showed that the adaptive selector was not the primary problem. It was
partially protecting the graph stage.

\subsection{Oracle-Node Test}

To distinguish proposal failure from architecture failure, the project added:

\begin{lstlisting}
scripts/diagnose_beast_oracle.py
\end{lstlisting}

This script replaces learned proposal nodes with GT-derived oracle nodes and
then runs the same foreground-geodesic initializer. If the graph stage works
with oracle nodes, the problem is proposal-side. If it still fails, the graph
architecture itself is insufficient.

Oracle result:

\begin{lstlisting}
oracle early hard samples:
  mean F1        0.5889
  mean bridges   76.6000
  mean reachable 0.3969
\end{lstlisting}

\begin{center}
\begin{tabular}{lccc}
\toprule
Sample & F1 & Bridges & Reachable \\
\midrule
7 & 0.2677 & 124 & 0.0236 \\
16 & 0.5433 & 104 & 0.1811 \\
22 & 0.8898 & 0 & 1.0000 \\
23 & 0.6373 & 73 & 0.4252 \\
24 & 0.6063 & 82 & 0.3543 \\
\bottomrule
\end{tabular}
\end{center}

Conclusion: most of these hard samples break PointNeuron1.0 even with good
nodes. The failure is architectural:

\begin{lstlisting}
saturated / huge / anisotropic volume
-> extreme foreground threshold adaptation
-> eligible geodesic pairs collapse
-> MST still forces one connected tree
-> fake long bridge edges
\end{lstlisting}

These samples should be treated as PointNeuron2.0 work, not as clean
PointNeuron1.0 training targets.

\reportbox{Architectural diagnosis}{
The early hard samples are not being avoided because they are inconvenient.
They are separated because the oracle-node test shows a real graph-stage
failure mode: the current PointNeuron1.0 geodesic tree builder can invent long
bridges even when the node candidates are already near the ground truth.
}

\section{Connectivity GAE Baseline}

The repository includes a PointNeuron-style graph autoencoder:

\begin{lstlisting}
src/pointneuron/models/connectivity.py
\end{lstlisting}

The GAE receives node features and an initial adjacency, then predicts an
adjacency matrix. In other words, it tries to learn which selected skeleton
nodes should be connected after seeing both node-level information and the
initial geodesic graph.

It was trained only on graph-eligible samples to avoid learning from corrupted
bridge-heavy targets. A sample was considered eligible when the initializer
already produced enough usable graph signal: sufficient nodes, acceptable edge
F1, acceptable reachable fraction, and no foreground-cap failure. A
foreground-cap failure means the foreground extraction produced a point set that
violated the practical limits used by the graph construction path, usually
because the sample's intensity/size behavior made the foreground too large or
unstable to treat as normal evidence. This was a
training-control decision. It prevented the first GAE baseline from being
dominated by cases where the target graph itself was already contaminated by
obvious reconstruction failure.

Eligibility result:

\begin{lstlisting}
accepted records: 46
rejected records: 39

rejection reasons:
  foreground_cap_not_satisfied: 5
  low_edge_f1: 17
  low_reachable: 13
  too_few_nodes: 4
\end{lstlisting}

Training result:

\begin{lstlisting}
records: 46
best_epoch: 46
best_loss: 0.377377
checkpoint: tmp/checkpoints/connectivity_guarded30_eligible_50e.pt
\end{lstlisting}

Evaluation result:

\begin{lstlisting}
threshold_precision mean 0.3702
threshold_recall    mean 0.9490
threshold_f1        mean 0.4983

topk_f1             mean 0.5200
init_topk_f1        mean 0.8261
\end{lstlisting}

The threshold metrics use a fixed probability cutoff of 0.5. The top-k metrics
instead select the highest-scoring predicted edges until the same number of
edges as the target graph is reached. The comparison to
\texttt{init\_topk\_f1} asks whether the learned GAE improves over the graph it
was given as input. In this run it did not.

Conclusion: the GAE learned a high-recall, low-precision edge signal. It did
not beat the initialized geodesic graph and is not promoted.

\reportbox{Connectivity decision}{
The learned GAE is useful as a baseline and a diagnostic artifact, but not as
the deployed PointNeuron1.0 connectivity stage. Its precision is too low, and
its top-k F1 remains below the initializer it was meant to improve.
}

\section{Final PointNeuron1.0 Status}

Promoted:

\begin{lstlisting}
proposal checkpoint:
  tmp/checkpoints/proposal_conservative_guarded30_cw3_oreg005_nw3.pt

topology path:
  adaptive_connected_coverage_nms geodesic graph
\end{lstlisting}

Not promoted:

\begin{lstlisting}
tmp/checkpoints/proposal_conservative_rawcenter_full60_cw3_oreg001_nw1.pt
tmp/checkpoints/connectivity_guarded30_eligible_50e.pt
\end{lstlisting}

Boundary:

\begin{lstlisting}
PointNeuron1.0 is suitable for the stable validated Gold166 subset.
PointNeuron1.0 is not architecture-complete for the early hard samples.
The first connectivity GAE does not replace the geodesic initializer.
\end{lstlisting}

\section{PointNeuron2.0 Implications}

The PointNeuron2.0 proposal is aligned with the observed failure, but the most
important fix is learned connectivity. Better proposals alone are not enough,
because the oracle-node experiment showed that even GT-derived nodes fail under
the current foreground-geodesic forced-tree rule.

Recommended PointNeuron2.0 direction:

\begin{enumerate}
  \item Use stronger point-cloud encoding and proposal generation, such as the
  proposed PointNeXt / transformer-decoder detection approach.
  \item Build multi-scale candidate edges rather than scoring all possible
  pairs equally.
  \item Train an edge classifier or reranker using learned features, Euclidean
  geometry, radius information, local direction, and foreground-geodesic
  evidence.
  \item Treat foreground geodesic as a feature, not the final source of truth.
  \item Apply topology constraints after edge scoring rather than forcing a
  tree before edge evidence is reliable.
\end{enumerate}

\newpage

\section{Reproducibility Commands}

\subsection{Environment}

\begin{lstlisting}[language=bash]
py scripts\check_environment.py --config configs\local.json
py -m pip install -e .
py -m pip install scipy
\end{lstlisting}

\subsection{Manifest and Inspection}

\begin{lstlisting}[language=bash]
py scripts\build_gold166_manifest.py `
  --root data\gold166 `
  --output tmp\gold166_manifest.json

py scripts\inspect_volume.py --sample-index 0 --decode

py scripts\check_sample_alignment.py `
  --sample-index 0 `
  --decode-volume
\end{lstlisting}

\subsection{Current Baseline}

\begin{lstlisting}[language=bash]
py scripts\run_end_to_end.py `
  --sample-index 110 `
  --output-root tmp\e2e_guarded30_single `
  --checkpoint tmp\checkpoints\proposal_conservative_guarded30_cw3_oreg005_nw3.pt `
  --initializer geodesic `
  --graph-selection-mode adaptive_connected_coverage_nms `
  --graph-min-proposal-score 0.85 `
  --graph-nms-mode distance `
  --graph-nms-distance 18 `
  --graph-max-nodes 128 `
  --device cuda
\end{lstlisting}

\subsection{Topology Evaluation}

\begin{lstlisting}[language=bash]
py scripts\evaluate_geodesic_baseline.py `
  --summary tmp\e2e_guarded30_adaptive_v4_110_129\summary.json `
  --csv-output tmp\e2e_guarded30_adaptive_v4_110_129\topology_report.csv `
  --json-output tmp\e2e_guarded30_adaptive_v4_110_129\topology_report.json
\end{lstlisting}

\subsection{Early Hard Oracle Diagnosis}

\begin{lstlisting}[language=bash]
py scripts\diagnose_beast_oracle.py `
  --output-root tmp\beast_oracle_diagnosis
\end{lstlisting}

\subsection{Connectivity Baseline Reproduction}

\begin{lstlisting}[language=bash]
py scripts\build_eligible_connectivity_manifest.py `
  --summary tmp\e2e_guarded30_adaptive_v4_000_024\summary.json `
  --summary tmp\e2e_guarded30_adaptive_v4_100_109\summary.json `
  --summary tmp\e2e_guarded30_adaptive_v4_110_129\summary.json `
  --summary tmp\e2e_guarded30_adaptive_v4_130_162\summary.json `
  --output-root tmp\connectivity_guarded30_eligible `
  --include-score `
  --skip-existing

py scripts\train_connectivity.py `
  --record-manifest tmp\connectivity_guarded30_eligible\eligible_connectivity_manifest.json `
  --epochs 50 `
  --lr 1e-3 `
  --weight-decay 1e-4 `
  --normalize-node-features `
  --init-decoder-bias `
  --checkpoint tmp\checkpoints\connectivity_guarded30_eligible_50e.pt `
  --device cuda
\end{lstlisting}

\section{Conclusion}

The project successfully reproduced the main PointNeuron1.0 pipeline stages and
produced a stable guarded30 + adaptive geodesic baseline for the validated
domain. It also clarified two important negative results. First, lower proposal
validation loss does not necessarily imply better reconstruction. Second, the
current full-adjacency connectivity GAE does not outperform the geodesic graph
initializer.

Most importantly, the oracle-node diagnosis established a real architectural
boundary. Several early hard samples fail even when the graph stage is given
GT-derived nodes. This means the problem is not merely insufficient proposal
training. It is a limitation of the current foreground-geodesic forced-tree
graph model. The next major research step should therefore be a PointNeuron2.0
connectivity architecture that learns candidate-edge reliability and treats
foreground geodesic evidence as one feature rather than as the final authority.

\end{document}
